%% LyX 2.4.4 created this file.  For more info, see https://www.lyx.org/.
%% Do not edit unless you really know what you are doing.
\documentclass[english]{article}
\usepackage[T1]{fontenc}
\usepackage[latin9]{inputenc}
\usepackage{float}
\usepackage{units}
\usepackage{enumitem}
\usepackage{amsmath}
\usepackage{amsthm}
\usepackage{graphicx}
\usepackage{geometry}
\geometry{verbose,tmargin=1in,bmargin=1in,lmargin=1in,rmargin=1in}

\makeatletter

%%%%%%%%%%%%%%%%%%%%%%%%%%%%%% LyX specific LaTeX commands.
%% A simple dot to overcome graphicx limitations
\newcommand{\lyxdot}{.}


%%%%%%%%%%%%%%%%%%%%%%%%%%%%%% Textclass specific LaTeX commands.
\numberwithin{equation}{section}
\numberwithin{figure}{section}
\newlength{\lyxlabelwidth}      % auxiliary length 

%%%%%%%%%%%%%%%%%%%%%%%%%%%%%% User specified LaTeX commands.
\usepackage{fancyhdr}
\usepackage{lastpage}
\pagestyle{fancy}
\usepackage{enumitem}
\usepackage{ifthen}
%sets math fonts to be more readable
\everymath=\expandafter{\the\everymath\displaystyle}
%\DeclareMathSizes{10}{10}{10}{10}
\usepackage{multicol}
% Added by lyx2lyx
\setlength{\parskip}{\smallskipamount}
\setlength{\parindent}{0pt}

\makeatother

\usepackage{babel}
\begin{document}
\renewcommand{\author}{Dr.\,James Ashe}

\newcommand{\classNum}{331}

\newcommand{\classSec}{A}

\newcommand{\nameType}{Name:}

\newcommand{\examType}{Exam 2}

\renewcommand{\date}{November 11, 2012}

%%First page's header%%%%%%%%%%%%%%%%%%%%%%
\fancypagestyle{firststyle} %%header settings for first pagr
{
	\fancyhead[L]{MTH \classNum{}(\classSec{}) \author{}\\
	\textbf{\examType{}} \date{}}
	\fancyhead[C]{\textbf{\nameType}}
	\setlength{\headsep}{14pt}
}
%%%%%%%%%%%%%%%%%%%%%%%%%%%%%%%%%%%%%%%%%%

%%Next page's header%%%%%%%%%%%%%%%%%%%%%%%
\setlist{nolistsep}
\lhead{\classNum{}(\classSec{})} 
\chead{\textbf{\nameType}} 
\cfoot{\thepage{}~of~\pageref{LastPage}} 
\setlength{\headsep}{5pt} 
%%%%%%%%%%%%%%%%%%%%%%%%%%%%%%%%%%%%%%%%%%%

%%Various settings; see the preamble for other settings%%%%%
%%\setenumerate[0]{leftmargin=12pt,itemindent=0pt,labelwidth=0pt}
\renewcommand{\labelenumi}{\textbf{\arabic{enumi}.}}
\renewcommand{\labelenumii}{\textbf{\alph{enumii}.}}
\thispagestyle{firststyle}
%%%%%%%%%%%%%%%%%%%%%%%%%%%%%%%%%%%%%%%%%%

\emph{Determine the radius of convergence of the following power series. }
\begin{enumerate}[resume]
\item ${\displaystyle \sum\frac{x^{k}}{2^{k}}}$\vfill
\end{enumerate}
\emph{Use the geometric series ${\displaystyle \sum_{k=0}^{\infty}x^{k}=\frac{1}{1-x}}$,
for $\left|x\right|<1$ to find a power series representations for
the following functions, and determine the interval of convergence.}
\begin{enumerate}[resume]
\item $f(x)=\frac{10}{1-x^{2}}$.\vfill
\item $f(x)=\frac{2x^{2}}{1-x}$.\vfill
\end{enumerate}
\emph{Eliminate the parameter in following equations to obtain a single
equation in terms of $x$ and $y$. }
\begin{enumerate}[resume]
\item $x=2t$, $y=4t^{2}$.\vfill
\end{enumerate}
\emph{Determine $\nicefrac{dy}{dx}$ in terms of $t$ and evaluate
it at the given value of $t$. }
\begin{enumerate}[resume]
\item $x=2t$, $y=4t^{2}$; $t=1$.\vfill
\end{enumerate}
\emph{Give a set of parametric equations that describes the curve.}
\begin{enumerate}[resume]
\item A circle centered at $(0,0$) with radius 6, generated clockwise.\vfill
\item \emph{Express the following }\textbf{\emph{polar}}\emph{ coordinates
in }\textbf{\emph{Cartesian}}\emph{ coordinates.}
\item $\left(-4,\frac{2\pi}{3}\right)$\vfill
\end{enumerate}
\emph{Express the following }\textbf{\emph{Cartesian}}\emph{ coordinates
in }\textbf{\emph{polar}}\emph{ coordinates in at least }\textbf{\emph{two}}\emph{
different ways. }
\begin{enumerate}[resume]
\item $\left(-1,1\right)$\vfill\newpage
\end{enumerate}
\emph{Locate the following polar coordinates on the graph. }
\begin{enumerate}[resume]
\item A:$\left(2,\frac{\pi}{4}\right)$, B:$\left(-2,\frac{\pi}{4}\right)$,
C:$\left(2,\frac{\pi}{4}+\pi\right)$, D:$\left(2,\frac{\pi}{4}+2\pi\right)$
\begin{figure}[H]
\begin{centering}
\includegraphics[width=6cm]{D:/home/jim/JamesAshe/JCSU/Calc_3_MTH_331/slides_Calc3/images/polar_graph_-2-2a.png}
\par\end{centering}
\caption{Use this graph to plot the polar coordinates.}
\end{figure}
\end{enumerate}
\emph{Mark the points on the }\textbf{\emph{polar}}\emph{ graph that
corresponds to the points shown on the }\textbf{\emph{Cartesian}}\emph{
graph.}
\begin{enumerate}[resume]
\item Mark the following points: A, B, C, D. 
\begin{figure}[H]
\begin{centering}
\emph{\includegraphics[width=6cm]{D:/home/jim/JamesAshe/JCSU/Calc_3_MTH_331/slides_Calc3/images/graph_2sin(t)-1.5.png}\includegraphics[width=6cm]{D:/home/jim/JamesAshe/JCSU/Calc_3_MTH_331/slides_Calc3/images/polar_graph_2sin(t)-1.5.png}}
\par\end{centering}
\caption{The Cartesian graph is on the left and the polar graph is on the right.}
\end{figure}
\end{enumerate}
\emph{Find the slope of the line tangent to the following polar curves
at the given point using calculus. }
\begin{enumerate}[resume]
\item $r=2\cos2\theta$; $\left(2,\frac{\pi}{2}\right)$.\vfill
\end{enumerate}
\emph{Find the area of the region using calculus. }
\begin{enumerate}[resume]
\item The region enclosed by all the leaves of the rose $r=2\sin4\theta$.\vfill
\end{enumerate}

\end{document}
